\heading{经济学原理（II）-课程说明}
\noteinfo{本课程说明依据你上传的第 1--6 讲讲义与样题风格整理，侧重考试中最常考、最容易混淆、最需要会“解释为什么”的部分。以下不是简单罗列名词，而是把每个概念放回其经济学含义中。}

\textbf{一、这门课在考什么}

《经济学原理（II）》表面上分散在税制、GDP、增长、金融与风险等章节，实际一直围绕同一个主线：\textbf{如何用一套统一的经济学语言理解“总量运行、长期增长、分配公平与风险配置”}。因此复习时不能把每一章看成孤立板块，而应当反复问自己四个问题：

\begin{itemize}[leftmargin=2em,itemsep=0.25em,topsep=0.25em]
  \item 这个变量描述的是\textbf{流量}还是\textbf{存量}？
  \item 这个结论是关于\textbf{水平}，还是关于\textbf{变化率}？
  \item 这个政策影响的是\textbf{效率}、\textbf{公平}，还是二者之间的权衡？
  \item 这个模型说的是\textbf{短期调整}，还是\textbf{长期趋势}？
\end{itemize}

考试最爱出的，不是单纯定义题，而是“你是否真的知道这些概念为什么要这样定义”。

\textbf{二、核心概念归纳（带解释）}

\textbf{1. 权衡取舍、机会成本与边际思维}

经济学首先不是算式，而是思维方式。资源稀缺意味着你不可能同时得到一切，于是任何选择都包含\textbf{权衡取舍}。机会成本不是“花了多少钱”，而是\textbf{为了得到某个选择而放弃的最佳备选项价值}。这一定义很重要，因为它告诉你：判断一个决策是否值得，关键不是看账面支出，而是看你因此失去了什么最好的替代选择。

边际思维则进一步告诉你，很多决策不是“做或不做”，而是“再多做一点是否值得”。宏观政策同样如此：财政扩张可以稳增长，但如果通胀压力已高，再多刺激一点的边际收益可能低于边际代价；减税可以增强激励，但减税过度又可能削弱财政能力。考试里只要看到“是否继续扩大”“再增加一点”“额外 1 单位”的表述，基本就在考边际分析。

\textbf{2. 宏观与微观的关系}

微观研究个体行为与单个市场，宏观研究总产出、总就业、总物价等总量变量。两者不是割裂的。宏观中的消费、储蓄、投资、劳动供给，本质上都来自微观个体行为的汇总；但一旦上升到总量层面，又会出现单个主体看不到的结果，例如失业、通胀、总需求不足和金融风险。考试中常见的陷阱，是把个体最优误认为总体最优；例如家庭都想多储蓄，在个体层面合理，但若总需求过弱，可能反而压低收入。

\textbf{3. 税收、公共品与再分配}

政府征税并不只是为了“有钱可花”，而是因为市场并不总能自动给出最有效、最公平的结果。公共品具有\textbf{非竞争性}和\textbf{非排他性}，私人市场往往供给不足；外部性意味着私人决策没有把全部社会成本或社会收益算进去；再分配则回应了市场结果可能带来的过大收入差距。因此税收既是\textbf{筹资工具}，也是\textbf{激励工具}与\textbf{分配工具}。

但税制设计永远面对效率与公平的张力。总额税对行为扭曲最小，所以效率高；但它不考虑支付能力，公平性弱。累进税更强调纵向公平，却可能改变劳动供给、储蓄与投资激励。考试中看到“边际税率”和“平均税率”时要立刻想到：\textbf{平均税率主要描述税负分担，边际税率主要影响行为激励}。

\textbf{4. 收入分配、不平等与贫困}

不平等不是简单地“有差距就是坏事”。适度差距可能反映人力资本投资、努力程度和承担风险的差异，从而形成激励；但若差距过大，可能损害社会流动、抑制消费并引发长期问题。衡量不平等时，洛伦兹曲线描述“人口累计占比”与“收入累计占比”的关系；离 45° 线越远，说明越不平等。基尼系数就是把这种“弯曲程度”压缩成一个数字。它不是凭空定义出来的，而是\textbf{把收入分配相对均等线偏离的面积比例化}。

贫困则与不平等不同：不平等讲的是“分配差距”，贫困讲的是“是否低于某种基本生活标准”。因此一种政策可以降低贫困，但不一定显著降低基尼系数；反过来，一种政策也可能压低高收入群体份额，却对最贫困群体帮助不大。

\textbf{5. GDP、GNP 与“生产”口径}

GDP 衡量的是\textbf{一国境内}在一定时期内生产的最终产品和服务的市场价值；GNP（或 GNI）衡量的是\textbf{一国居民}所获得的生产收入。二者差别的本质不是“哪个更先进”，而是一个强调\textbf{地域口径}，一个强调\textbf{居民口径}。

为什么 GDP 只算最终产品？因为如果把面粉、面包、面包店销售额都加一遍，就会重复计算。同理，旧房、二手车本身不计入当期 GDP，因为它们不是本期生产；但中介服务、经纪佣金会计入 GDP，因为那是本期新增服务。转移支付不计入 G，也不是因为它“不重要”，而是因为它不是对本期生产的支付；它会通过改变居民可支配收入，\textbf{间接}影响消费和 GDP。

\textbf{6. 名义变量、真实变量与价格指数}

名义 GDP 用当期价格计价，真实 GDP 用基期价格或链式方法剔除价格变化影响。二者的差别告诉我们：经济总量变大，究竟是因为\textbf{产出真的更多了}，还是因为\textbf{价格更贵了}。GDP 平减指数反映本国生产品的总体价格；CPI 则反映代表性消费者消费篮子的生活费用变化。两者不同，是因为覆盖对象不同：GDP 平减指数关注“国内生产了什么”，CPI 关注“居民消费了什么”。

CPI 常被认为会高估生活费用上升，其中最核心的原因是\textbf{替代偏差}：当某种商品涨价时，消费者会减少它、转向替代品，而固定篮子指数没有反映这种行为调整。考试中只要看到“固定篮子”“生活费用”“代表性消费者”，多半就是 CPI；看到“国内生产”“当期产出构成”，通常是 GDP 平减指数。

\textbf{7. 经济增长、70 规则与长期决定因素}

增长率真正重要，是因为它反映福利与国力的变化速度。70 规则不是魔法公式，而是对低增长率下指数增长的近似：若某变量按每年 $g\%$ 增长，则翻倍所需时间约为 $70/g$ 年。它帮助我们快速把“抽象增长率”转化成“需要多久翻倍”的直观判断。

长期增长来自劳动、资本、人力资本、自然资源、技术进步和制度安排，但在现代增长理论里，真正决定长期持续增长能力的关键是\textbf{技术进步与效率提升}。单纯依赖资本积累会遭遇边际报酬递减；资本可以推动从低水平向稳态收敛，却难以单独支撑永久高增长。

\textbf{8. 索罗模型中的稳态、资本深化与黄金规则}

索罗模型最核心的问题是：\textbf{每个劳动者能配到多少资本}。记人均资本为 $k$，人均产出为 $y=f(k)$。储蓄把一部分产出转化为投资，推动资本深化；而人口增长与折旧会“摊薄”资本，形成资本广化压力。稳态不是“不增长”，而是\textbf{人均资本不再变化}，也就是“实际人均投资”正好等于“维持资本密度不变所需投资”。

黄金规则储蓄率追求的不是最大产出，而是\textbf{最大稳态人均消费}。这点极易混淆：一个社会把全部产出都拿去投资，未必让居民过得最好；最优长期状态是让“今天少消费一些带来的未来增产”与“今天少消费本身的代价”达到平衡。

\textbf{9. 储蓄、投资与可贷资金市场}

在宏观上，储蓄不是单纯的“美德”，而是为投资提供资源；投资则决定资本积累和未来生产能力。可贷资金市场用一个极简框架刻画：真实利率调节储蓄供给与投资需求。当政府出现预算赤字时，公共储蓄下降，国民储蓄减少，可贷资金供给左移，真实利率上升，投资被部分挤出；当预算盈余扩大时，方向相反。

因此财政政策不仅影响总需求，也影响资本形成与长期增长路径。考试里只要看到“预算赤字”“可贷资金”“挤出效应”，就要立刻画出储蓄供给左移、真实利率上升、投资下降的链条。

\textbf{10. 名义利率、真实利率与费雪效应}

名义利率是合同写出来的利率，真实利率反映购买力层面的回报。两者关系近似为
\[
i \approx r + \pi^e.
\]
这不是单纯记忆题。它意味着：如果预期通胀上升，而真实回报要求不变，那么名义利率会上升。换言之，货币现象会“嵌入”名义利率；真正引导储蓄与投资决策的，通常是\textbf{真实利率}。

\textbf{11. 风险、收益与分散化}

金融学里“风险高收益高”不是道德判断，而是说高收益往往对应更大的不确定性。风险可以用方差或标准差刻画；期望收益衡量平均回报。分散化的经济学含义是：如果不同资产的波动不完全同步，把它们放进组合后，个别资产的特殊波动可以相互抵消，所以能降低\textbf{非系统性风险}。但整个经济周期、利率冲击、金融危机这类风险无法仅靠多买几只股票消掉，它们属于\textbf{系统性风险}。

\textbf{12. 有效市场假说的正确理解}

有效市场假说并不是“市场永远正确”，而是说\textbf{公开信息会被迅速反映到价格里}。它强调的是“持续击败市场并不容易”，不是说价格从不偏离基本面。现实中，信息处理、行为偏差、制度摩擦都可能使价格阶段性偏离，因此有效市场与资产泡沫并不是一个简单的“二选一”。

\textbf{三、核心公式整理（带意义解释）}

\textbf{1. GDP 支出法}
\[
Y=C+I+G+NX.
\]
这里的 $Y$ 是总产出，也是总收入；$C$ 是消费，$I$ 是投资，$G$ 是政府购买，$NX=X-M$ 是净出口。最常见错误有三个：把转移支付算进 $G$；把新房算进 $C$；忘记进口品会先进入 $C/I/G$ 再在 $NX$ 里扣掉。

\textbf{2. GDP 增长率与人均 GDP 增长率}
\[
g_Y=\left(\frac{Y_t}{Y_{t-1}}-1\right)\times 100\%,\qquad
g_y=\left(\frac{y_t}{y_{t-1}}-1\right)\times 100\%.
\]
其中 $y=Y/N$。在增长率不高时有近似关系
\[
g_y \approx g_Y-g_N.
\]
这说明总量增长未必等于人均福利改善；若人口增长快，人均增长会被“吃掉”一部分。

\textbf{3. 70 规则}
\[
\text{翻倍时间} \approx \frac{70}{g(\%)}.
\]
它是指数增长近似，不是精确恒等式，但在考试中非常常用。

\textbf{4. 消费函数、MPC 与 MPS}
\[
C=C_0+\alpha(Y-T),\qquad MPC=\frac{\Delta C}{\Delta(Y-T)},\qquad MPS=\frac{\Delta S}{\Delta(Y-T)}.
\]
并且
\[
MPC+MPS=1.
\]
$C_0$ 表示即使没有收入也要发生的基本消费；$\alpha$ 表示可支配收入每增加 1 单位，有多少会转为消费。MPC 越大，居民更愿意把新增收入花掉；MPS 越大，则更愿意储蓄。

\textbf{5. 国民储蓄恒等式}
在封闭经济中，
\[
Y=C+I+G,\qquad S=Y-C-G,\qquad S=I.
\]
它不是说“现实中任何时刻储蓄现金一定等于投资现金”，而是国民收入核算下的宏观恒等关系。进一步可分解为
\[
S=(Y-T-C)+(T-G),
\]
即\textbf{私人储蓄 + 公共储蓄}。

\textbf{6. 费雪方程}
\[
i \approx r+\pi^e.
\]
名义利率 $i$ 体现合同回报，真实利率 $r$ 体现购买力回报，$\pi^e$ 是预期通胀。该公式解释了为什么高通胀预期时代，名义利率通常也高。

\textbf{7. 不平等度量：基尼系数}
若以五等份数据近似，基尼系数可由洛伦兹曲线下方面积 $A_L$ 得到：
\[
G=1-2A_L.
\]
本质上它衡量“实际分配曲线”偏离绝对平等线的程度。用五等份或十分位数据时，本质上是在用折线近似真实曲线，因此是近似值。

\textbf{8. 索罗模型的核心方程}
设 $y=f(k)$，储蓄率为 $s$，人口增长率为 $n$，折旧率为 $\delta$，则稳态条件是
\[
sf(k)=(n+\delta)k.
\]
左边是实际人均投资，右边是维持资本密度不变所需投资。若左边更大，$k$ 上升；若右边更大，$k$ 下降。

\textbf{9. 黄金规则}
黄金规则的思想是最大化稳态消费
\[
c^*=f(k^*)-(n+\delta)k^*.
\]
若生产函数是 Cobb--Douglas 形式 $y=k^\alpha$，则黄金规则储蓄率常等于
\[
s_{\mathrm{gold}}=\alpha.
\]
这不是因为“资本份额神奇地等于最优”，而是因为在该生产函数下，边际产出递减的具体形式恰好导致这一结果。

\textbf{10. 期望收益与风险}
若资产在状态 $s$ 下收益率为 $r_s$，概率为 $p_s$，则
\[
E(r)=\sum_s p_s r_s,\qquad
\mathrm{Var}(r)=\sum_s p_s(r_s-E(r))^2.
\]
标准差是方差的平方根。期望收益反映平均回报，标准差反映波动性。组合理论的核心不在于“算数难”，而在于理解相关性如何影响总风险。

\textbf{四、最容易丢分的地方}

\begin{itemize}[leftmargin=2em,itemsep=0.25em,topsep=0.25em]
  \item 把\textbf{新房}误记为消费；实际上属于投资。
  \item 把\textbf{转移支付}算进政府购买；实际上它不对应本期生产。
  \item 只会算 GDP 总量，不会区分\textbf{名义增长}和\textbf{真实增长}。
  \item 看到预算赤字，只会背“政府没钱了”，却不会写出\textbf{公共储蓄下降 $\to$ 国民储蓄下降 $\to$ 真实利率上升 $\to$ 投资下降}。
  \item 在索罗模型里只会背公式，不会比较 $sf(k)$ 与 $(n+\delta)k$ 的大小。
  \item 把\textbf{平均税率}和\textbf{边际税率}混淆。
  \item 知道“分散化降低风险”，但说不清它主要降低的是\textbf{非系统性风险}。
\end{itemize}

\textbf{五、如何使用后面的两套模拟题}

这两套模拟题刻意把难度抬高到“比样题更综合一步”的水平，主要体现在：概念之间会交叉、计算题需要多步解释、部分选择题会考错误选项为什么错。建议做法是：

\begin{itemize}[leftmargin=2em,itemsep=0.25em,topsep=0.25em]
  \item 第一遍严格限时，把不会的先跳过，训练考试节奏；
  \item 第二遍对照答案，不只看“对错”，而要总结自己是\textbf{概念错、口径错、公式错还是逻辑链条断了}；
  \item 对每道计算题，都试着用一句话概括它在考什么，例如“这题本质在考 GDP 口径判断”或“这题本质在考稳态条件而不是机械代数”；
  \item 如果时间有限，优先反复刷：GDP 核算、价格指数、基尼系数、索罗稳态、可贷资金市场、风险分散。
\end{itemize}

\textbf{六、总体复习建议}

这门课真正拉开差距的，不是背得多，而是\textbf{能不能把定义翻译成经济含义}。如果你能做到以下三点，分数通常就会比较稳：

\begin{itemize}[leftmargin=2em,itemsep=0.25em,topsep=0.25em]
  \item 一看到题目，就先判断它属于哪个“母题”；
  \item 一写公式，就知道每个符号背后的经济含义；
  \item 一下结论，就能补上一句“为什么会这样”。
\end{itemize}

做到这三点，模拟题就不只是刷题，而会真正变成建立框架的过程。

